%% ============================================================
%%  Section: Methodology
%%  Skeleton — \placeholder{} marks where formal derivations are needed.
%% ============================================================

\section{Methodology}
\label{sec:method}

\subsection{Environments}
\label{sec:method:envs}

We evaluate on two visual RL benchmarks:

\paragraph{Visual Pursuit.}
A smooth continuous-control tracking task in which an agent controls a
reticle $(x, y)$ to track a smoothly-moving target on a
$64 \times 64 \times 3$ observation canvas.
Perturbations include photometric shifts (brightness, contrast, hue) and
camera noise.
Reward is the negative Euclidean tracking error $r = -\|p_\mathrm{agent} -
p_\mathrm{target}\|_2$.

\paragraph{Procgen CoinRun.}
A discrete combinatorial game environment from the Procgen
benchmark~\citep{\placeholder{cite_procgen}}.
Four perturbation axes are evaluated: \textit{background}, \textit{sprite},
\textit{wall texture}, and \textit{difficulty (hazard)}.

\subsection{Classical Representation Distance \texorpdfstring{$d_\mathrm{rep}$}{d\_rep}}
\label{sec:method:drep}

Given a frozen encoder $\phi$, the representation distance between
observation distribution $\mathcal{D}_\mathrm{clean}$ and perturbed
distribution $\mathcal{D}_\mathrm{pert}$ is:
\[
  d_\mathrm{rep}(\mathcal{D}_\mathrm{clean}, \mathcal{D}_\mathrm{pert})
  = \placeholder{formal definition of d\_rep}
\]

\subsection{Semantically-Weighted Distance \texorpdfstring{$\tilde{d}_\mathrm{rep}$}{d-tilde\_rep}}
\label{sec:method:dtilderep}

To address metric collapse on combinatorial shifts, we introduce
weighter functions $w: \mathcal{X} \to [0, 1]$ that mask irrelevant
background pixels:
\[
  \tilde{d}_\mathrm{rep}(\cdot)
  = \placeholder{formal definition of d-tilde\_rep with weighter w}
\]

We evaluate four weighter backends:
(i) \emph{Classical} (no weighting),
(ii) \emph{SegmentationWeighter} using ground-truth segmentation masks,
(iii) \emph{TLDAWeighter} using Lipschitz-estimated policy sensitivity, and
(iv) \emph{SaliencyWeighter} using critic gradient saliency maps.

\subsection{Grassmann Subspace Geometry}
\label{sec:method:grassmann}

We extract the top-$k$ principal subspace $U \in \mathrm{Gr}(k, D)$
from encoder feature covariance matrices using truncated SVD.
The geodesic distance on the Grassmann manifold is:
\[
  d_\mathrm{geo}(U, V) = \|\Theta\|_2
\]
where $\Theta = (\theta_1, \ldots, \theta_k)$ are the $k$ principal angles
between subspaces $U$ and $V$.

\subsection{Test-Time Self-Supervised Adaptation (TTSS)}
\label{sec:method:ttss}

At test time we maintain a rolling window of $k = \placeholder{window\_size}$
recent observations and perform one step of self-supervised adaptation
minimizing:
\[
  \mathcal{L}_\mathrm{TTSS} = \mathcal{L}_\mathrm{smooth} +
  \lambda \mathcal{L}_\mathrm{moment}
\]
where $\mathcal{L}_\mathrm{smooth}$ is temporal frame-embedding consistency
and $\mathcal{L}_\mathrm{moment}$ is TTT++ moment matching~\citep{\placeholder{cite_ttt_plus_plus}}
pulling test statistics to stored offline training statistics.

\paragraph{Trust-Region Safeguard.}
An adaptation step is \emph{accepted} if and only if
$\mathcal{L}_\mathrm{TTSS}(\theta') < \mathcal{L}_\mathrm{TTSS}(\theta)$,
i.e.\ the candidate update strictly reduces the self-supervised loss.
Otherwise the original encoder weights are retained.
\placeholder{Add pseudocode Algorithm box for TTSS with trust-region.}

%% AUDIT NOTE: trust_region_acceptance_rate = 1.0 in the Phase 3 eval run.
%% The safeguard was never triggered empirically. Reframe claims accordingly
%% or re-run with more adversarial settings. See results/audit_report.json.
